\begin{abstract}
Model merging has emerged as a highly cost-effective paradigm for fusing task-specific expert weights into a single unified multi-task network without additional training.
However, standard merging methods operate in a static, zero-temperature Euclidean landscape, forcing straight-line parameter interpolation across highly non-convex loss boundaries.
This causes severe destructive interference and representation collapse, a phenomenon we characterize as ``parameter-space system frustration.''
To resolve this, we introduce \textbf{ThermoMerge} (Thermodynamic Model Merging), a paradigm-shifting framework that reformulates model merging through the lens of statistical mechanics and thermodynamics.
Instead of rigid parameter averaging, we thermalize model outputs by mapping task expert predictions to state probabilities within a finite-temperature canonical Boltzmann ensemble.
We introduce a \textbf{Thermodynamic Annealing Schedule (TAS)} to actively flatten non-convex optimization barriers during unsupervised test-time adaptation on sequential unlabeled streaming calibration streams.
On these streams, we optimize layer-wise merging coefficients and task-specific local temperatures (parameterized as task-wise thermal capacities in logit space) by minimizing a novel \textbf{Helmholtz Free Energy Discrepancy (F-Min)} objective, which serves as a natural, physically grounded regularizer.
Empirically, ThermoMerge mitigates the ``Overfitting-Optimizer Paradox'' that causes standard unregularized test-time adaptive merging methods to collapse.
Under a realistic sequential streaming test-time adaptation setting on a pre-trained ResNet-18 backbone, ThermoMerge achieves an outstanding multi-task average accuracy of \textbf{29.05\%}, outperforming static Task Arithmetic (\textbf{27.25\%}) and AdaMerging (\textbf{26.10\%}), and outperforming the highly competitive SyMerge baseline (\textbf{27.90\%}), establishing a novel, rigorous connection between deep learning optimization and statistical physics.
\end{abstract}
